%% Q1 ctx
\textbf{Read the following text and answer the Question Nos.~1 to 6.}\par\smallskip
The algorithmic theory of beauty takes the subjectivity of the observer into account and postulates that among several observations classified as comparable by a given subjective observer, the aesthetically most pleasing one is the one with the shortest description. This is closely related to the principles of algorithmic information theory and minimum description length. For example, mathematicians welcome simple proof with a short description in their formal language. Another concrete example describes an aesthetically pleasing human face whose proportions can be described by very few bits of information, drawing inspiration from 15th century proportion studies by Leonardo da Vinci. Going further, this theory explicitly distinguishes between what is beautiful and what is interesting. Interestingness corresponds to the first derivative of subjectively perceived beauty. Here, the premise is that any observer continually tries to improve the predictability of the observations by discovering regularities such as repetition and symmetry.

%% Q1 q
According to the algorithmic theory:

%% Q1
The theory says the most pleasing observation is the one with the \textbf{shortest description}. Beauty therefore increases as description length decreases: there is a \textbf{positive correlation between beauty and minimum description}.

%% Q2
Mathematicians ``welcome simple proof with a short description in their formal language''. The formal language of mathematics therefore values \textbf{simplicity and brevity}.

%% Q3
The passage says the face example draws ``inspiration from 15th century \textbf{proportion studies} by Leonardo da Vinci''. His studies \emph{support} the theory; he did not predict it.

%% Q4
The official key marks (A) as the statement not supported. The passage says only that the theory \emph{distinguishes} the interesting from the beautiful, with interestingness being the \emph{rate of change} (first derivative) of perceived beauty. It does not claim that what is interesting is not beautiful.

%% Q5
The observer improves predictability ``by discovering regularities such as \textbf{repetition and symmetry}'', i.e.\ by noticing patterns of similarity.

%% Q6
The passage is entirely about the \textbf{algorithmic theory of beauty}.

%% Q7
The Bhopal gas tragedy (2--3 December 1984, Union Carbide plant) released \textbf{methyl isocyanate (MIC)}, killing thousands. It is the world's worst industrial disaster.

%% Q8
The human audible range is about \textbf{20 Hz to 20 kHz}. Below it is infrasound; above it is ultrasound.

%% Q9
Contaminated drinking water spreads water-borne diseases such as \textbf{typhoid}, cholera and hepatitis A.
\why{Scurvy is vitamin-C deficiency; malaria is mosquito-borne; anaemia is a nutritional deficiency.}

%% Q10
In fresh-water lakes \textbf{phosphorus} is usually the limiting nutrient for algal (primary) production. Excess phosphate from detergents and fertilisers causes eutrophication.

%% Q11
CO$_2$, methane and water vapour are greenhouse gases. \textbf{Calcium carbonate} is a solid (limestone), not a gas.

%% Q12
A tsunami is caused by the \textbf{sudden displacement (motion) of the ocean floor}, usually by an undersea earthquake, landslide or volcanic eruption.

%% Q13
\textbf{FTP} is \textbf{File Transfer Protocol}, used to transfer files between computers over a network.

%% Q14
Anti-virus, authorised software and scanning external devices are recommended. \textbf{``Never connect to the internet''} is not a recommendation; the internet is essential, and safe use is what matters.

%% Q15
A bare string of bits such as 10110110, without context or processing, is \textbf{data}.
\tip{DIKW hierarchy: Data $\rightarrow$ Information $\rightarrow$ Knowledge $\rightarrow$ Wisdom.}

%% Q16
Information is meaningful, relevant, clear and \textbf{timely}. Its value \emph{changes with time}, so ``unchangeable with time'' is not a characteristic of information.

%% Q17
The choice of ICT tools depends on the available infrastructure, learners' access and support staff, i.e.\ \textbf{all of the above}.

%% Q18
A \textbf{soft copy} is displayed electronically, not printed. The \textbf{VDU (monitor)} is a soft-copy output device; a printer gives hard copy; a CD-ROM is a storage medium.

%% Q19
Scientific method: \textbf{observation $\rightarrow$ hypothesis $\rightarrow$ deduction of consequences $\rightarrow$ verification}.

%% Q20
Research by a practitioner (teacher, manager) to improve his or her own practice is \textbf{action research}.

%% Q21
When every unit has an equal (known) chance of selection, the procedure is \textbf{probability sampling}.

%% Q22
\textbf{Reliability} is the consistency of measurement: the same results on \textbf{replication} (repeated measurement).

%% Q23
\textbf{Philosophical research} relies on reasoning and conceptual analysis rather than empirical data, so it is non-empirical.

%% Q24
\textbf{Delimitations} are the boundaries the researcher deliberately sets on the study (population, variables, area, period), i.e.\ \textbf{setting limits to the problem}. \emph{Limitations} are constraints beyond the researcher's control.

%% Q25
Classroom teaching aims to help students \textbf{apply disciplinary knowledge to life}, not merely to pass examinations or compete.

%% Q26
Per the key, \textbf{study habits} are the learner factor with the potential to influence learning directly.

%% Q27
Effective teaching requires \textbf{knowledge of teaching methods and of learners} (pedagogy and psychology), along with content.

%% Q28
The lecture method's main advantage is that it is \textbf{suitable for teaching large groups} economically.

%% Q29
Teaching aids are used to reach \textbf{the students' level of understanding}, making ideas concrete and clear.

%% Q30
Internal assessment marks should be \textbf{revealed to students} so that they get feedback and the process stays transparent.

%% Q31
The effectiveness of communication is judged by \textbf{feedback}, which shows whether the message was understood as intended.

%% Q32
Language, facial expressions and acquaintance all influence face-to-face communication. The key treats \textbf{location} as having no influence.

%% Q33
Vertical, horizontal and crosswise (diagonal) are the \emph{directions} of communication flow in an organisation. The key gives \textbf{verbal communication} as the answer: it is classified by \emph{mode}, not by direction, so it does not belong to this set of types.

%% Q34
Per the key, noise \textbf{adds on to the message codes}: unwanted signals get added to the message and distort it.

%% Q35
The first requirement for classroom communication is \textbf{drawing the students' attention}. Without attention no message is received.

%% Q36
Answering questions, taking notes and cross-questioning the teacher are responses that feed back to the teacher. \textbf{Creating disturbance} is not a feedback loop in classroom communication.

%% Q37
Student representation in governing bodies helps \textbf{promote democratic citizenship}: participation, responsibility and a voice in decisions.

%% Q38
Environmental education draws on biology, chemistry, geography, economics and sociology, so it is \textbf{multi-disciplinary} research.

%% Q39
Distance education aims to \textbf{reach the unreached}: people who cannot attend regular institutions because of distance, work or circumstances.

%% Q40
Teaching in higher education is a vocation and mission: it is \textbf{not merely a job, though salaried}.

%% Q41
Values are best taught by example: the teacher must \textbf{behave ideally} (``values are caught, not taught'').

%% Q42
Accrediting bodies (NAAC, NBA) are \textbf{statutory or autonomous bodies to maintain and assure the quality of education}.

%% Q43 q
Complete the sequence of squares by choosing the correct figure from the given alternatives:\\[3pt]
\centering
\newcommand{\sq}[1]{\framebox[1.1cm][c]{\rule{0pt}{0.8cm}\parbox[c]{1cm}{\centering\tiny #1}}}
\sq{$\angle$}\,\sq{$\triangle\angle$}\,\sq{$\triangle\triangle\triangle$}\,\sq{$\triangle\triangle\triangle$\\$\triangle\angle$}\,\sq{$\triangle\triangle\triangle$\\$\triangle\triangle\triangle$\\$\angle$}\,\sq{?}

%% Q43 opts
\framebox{\parbox{0.9cm}{\centering\tiny$\triangle\triangle\triangle$\\$\triangle\triangle\triangle$\\$\triangle\triangle\triangle$}}
\framebox{\parbox{0.9cm}{\centering\tiny$\triangle\triangle\triangle$\\$\triangle\triangle\triangle$\\$\triangle\triangle\triangle$\\$\triangle\angle$}}
\framebox{\parbox{0.9cm}{\centering\tiny$\triangle\triangle\triangle$\\$\triangle\triangle\triangle$\\$\triangle\triangle\triangle$\\$\triangle\triangle\angle$}}
\framebox{\parbox{0.9cm}{\centering\tiny$\triangle\triangle\triangle$\\$\triangle\triangle\triangle$\\$\triangle\angle\angle$}}

%% Q43
The number of complete triangles grows $0, 1, 3, 4, 6, \ldots$ and the figure alternates between having a half-drawn triangle ($\angle$) and not having one. In the 3rd figure the angle is ``completed'' into a full row of triangles; likewise the next step completes the figure into \textbf{three full rows, i.e.\ nine triangles}, as in option (A).

%% Q44 q
Put the correct letter for `?' in the following diagram:\\[3pt]
\centering\begin{tabular}{ccc}
\framebox{C} & $\triangleright$ B & \textcircled{F}\\[2pt]
\framebox{D} & $\triangleright$ C & \textcircled{?}\\[2pt]
\framebox{E} & $\triangleright$ D & \textcircled{T}\\
\end{tabular}

%% Q44
The circle letter's position is the \textbf{product} of the other two positions:
\[\text{C}(3)\times\text{B}(2)=6=\text{F},\qquad \text{E}(5)\times\text{D}(4)=20=\text{T}\]
\[\text{D}(4)\times\text{C}(3)=12=\mathbf{L}\]

%% Q45 q
Take Universal set U $=\{$A, B, C, \ldots, Z$\}$ of all English alphabets. If set A contains the letters from GUJARAT, the members of set B are the alphabets from PUNJAB, the set C contains the alphabets from SIKKIM and the elements of the set D are the alphabets from TAMILNADU, then the Venn diagram for the sets A, B, C, D will be:

%% Q45 opts
\tikz[baseline=-0.5ex,scale=0.25]{\draw (-2,-2) rectangle (2,2); \draw (0,0.8) circle(0.8); \draw (-0.8,0) circle(0.8); \draw (0.8,0) circle(0.8); \draw (0,-0.8) circle(0.8);}
\tikz[baseline=-0.5ex,scale=0.25]{\draw (-2,-2) rectangle (2,2); \draw (-0.6,1.1) circle(0.7); \draw (-0.6,0) circle(0.7); \draw (-0.6,-1.1) circle(0.7); \draw (0.7,0.1) circle(0.7);}
\tikz[baseline=-0.5ex,scale=0.25]{\draw (-2,-2) rectangle (2,2); \draw (-0.9,-0.5) circle(0.7); \draw (0.1,0) circle(0.7); \draw (0.8,0.6) circle(0.7); \draw (0.7,-0.2) circle(0.7);}
\tikz[baseline=-0.5ex,scale=0.25]{\draw (-2,-2) rectangle (2,2); \draw (0,0.8) circle(0.9); \draw (-0.2,1.0) circle(0.35); \draw (0,-0.8) circle(0.9); \draw (-0.2,-1.1) circle(0.35);}

%% Q45
Write out the sets:
\begin{center}\small
\begin{tabular}{ll}
A (GUJARAT) & \{G, U, J, A, R, T\}\\ B (PUNJAB) & \{P, U, N, J, A, B\}\\ C (SIKKIM) & \{S, I, K, M\}\\ D (TAMILNADU) & \{T, A, M, I, L, N, D, U\}\\
\end{tabular}
\end{center}
A$\cap$B $=$ \{U, J, A\}, A$\cap$D $=$ \{T, A, U\}, B$\cap$D $=$ \{U, N, A\}, and A$\cap$B$\cap$D $=$ \{U, A\}. C shares \{I, M\} with D only and nothing with A or B.
So A, B, D overlap mutually and C overlaps only D, as in diagram \textbf{(C)}.

%% Q46
T is the middle tank and T1 is the tank immediately before it. From T1 the pilot counts 9 more tanks to reach the last one, so there are $9-1=8$ tanks after T. Since T is in the middle, there are also 8 tanks before T.
Total $=8+1+8=\mathbf{17}$.

%% Q47 q
Fill in the cell of `?' by putting the correct letter:\\[3pt]
\centering\begin{tabular}{|c|c|c|}\hline G & B & C\\\hline H & F & H\\\hline A & D & ?\\\hline\end{tabular}

%% Q47
Per the key the answer is \textbf{E}. With letter positions, the first and third columns then have the same total, $G+H+A=7+8+1=16=C+H+E=3+8+5$, balancing about the middle column.

%% Q48
Ox, giraffe, camel, goat, elephant and deer are \textbf{herbivores}; the \textbf{dog} is a carnivore (and a pet), so it is the odd one out.

%% Q49 q
Put the number in the empty box by choosing the correct digit from the given alternatives:\\[3pt]
\centering\begin{tabular}{|c|c|c|}\hline 7 & 3 & 6\\\hline 5 & 4 & 7\\\hline 8 & ? & 2\\\hline\end{tabular}

%% Q49
Each row adds up to 16: $7+3+6=16$ and $5+4+7=16$, so $8+?+2=16$ gives $?=\mathbf{6}$.

%% Q50 q
If $a*b=\dfrac{b}{a}$, then which of the following statements is \emph{not} true? ($a$ and $b$ are real numbers)

%% Q50 opts
$a*0=0,\ a\neq0$
$0*b=0,\ b\neq0$
$1*1=1$
$a*a=1,\ a\neq0$

%% Q50
$0*b=\dfrac{b}{0}$ is \textbf{undefined} (division by zero), so (B) is not true.
Check the others: $a*0=0/a=0$ \checkmark; $1*1=1$ \checkmark; $a*a=a/a=1$ \checkmark.

%% Q51 q
For which set of values of $p, q, P, Q$ and for which arithmetic operation $*$, the expression $\left(\dfrac{p}{P}\right)*\left(\dfrac{q}{Q}\right)=-1$?

%% Q51
Option (A): $\dfrac{-1}{-4}+\dfrac{5}{-4}=\dfrac14-\dfrac54=-1$ \checkmark.
(B) gives $1\times1=1$; (C) gives $-1\div-1=1$; (D) gives $1-(-1)=2$.

%% Q52
Two interleaved series: $12, 16, 20, 24$ ($+4$) and $15, 20, 25, \mathbf{30}$ ($+5$).

%% Q53
Positions double: A(1), B(2), D(4), H(8), \textbf{P(16)}.

%% Q54 q
If $l=\sqrt3$, $m=1/7$, $n=2$ and $t=5/4$, then $l^2 * m \;\square\; n \;o\; t=-11/28$ for

%% Q54 opts
$*$, $\square$ and $o$ are all `$-$'
$*$ is $+$, $\square$ is $+$, $o$ is `$-$'
$*$ is $-$, $\square$ is $-$, $o$ is `$+$'
$*$, $\square$ and $o$ are all `$+$'

%% Q54
With all three operations as minus: $3-\dfrac17-2-\dfrac54=\dfrac{84-4-56-35}{28}=-\dfrac{11}{28}$ \checkmark.

%% Q55 ctx
Answer the questions 55, 56 and 57 from the given data in graphical form. Given data of number of washing machines in thousands, manufactured by two companies A and B, during 1997 to 2002.
\begin{center}
\begin{tikzpicture}
\begin{axis}[width=10cm,height=5cm,xmin=1996.5,xmax=2002.5,ymin=40,ymax=175,xtick={1997,...,2002},
  x tick label style={/pgf/number format/1000 sep=},tick label style={font=\scriptsize},ylabel={\scriptsize thousands},
  nodes near coords,every node near coord/.append style={font=\tiny},legend style={font=\tiny,at={(1.02,1)},anchor=north west}]
\addplot[mark=square,thick,brand] coordinates {(1997,139) (1998,120) (1999,100) (2000,128) (2001,107) (2002,148)};
\addplot[mark=o,thick,dashed,accent] coordinates {(1997,119) (1998,99) (1999,141) (2000,78) (2001,120) (2002,159)};
\legend{A,B}
\end{axis}
\end{tikzpicture}
\end{center}

%% Q55 q
What is the difference between the number of washing machines manufactured by company B in 1997 and 1998?

%% Q55
B: $119-99=20$ thousand $=\mathbf{20{,}000}$.

%% Q56 q
What is approximate average number of washing machines manufactured by company A over the given period?

%% Q56
A: $\dfrac{139+120+100+128+107+148}{6}=\dfrac{742}{6}\approx123.67$ thousand $\approx\mathbf{123666}$.

%% Q57 q
In which year the difference between the production of companies A and B was minimum?

%% Q57
Differences (thousand): 1997: 20; 1998: 21; 1999: 41; 2000: 50; 2001: 13; 2002: \textbf{11}. The minimum is in \textbf{2002}.

%% Q58 ctx
Given the table of marks obtained by two candidates P and Q in four papers. Answer the following questions 58, 59 and 60 from this data.
\begin{center}\small
\begin{tabular}{lcccc}\toprule \textbf{Paper} & I & II & III & IV\\\midrule
Maximum marks & 80 & 70 & 100 & 60\\ Candidate P & 48 & 56 & 75 & 42\\ Candidate Q & 40 & 63 & 82 & 36\\\bottomrule
\end{tabular}
\end{center}

%% Q58 q
Arrange in increasing order the \% marks obtained by the candidate P. The sequence of papers is:

%% Q58
P's percentages: I $=48/80=60\%$, II $=56/70=80\%$, III $=75\%$, IV $=42/60=70\%$.
Increasing order: \textbf{I (60), IV (70), III (75), II (80)}.

%% Q59 q
In which paper(s) the candidate Q obtained the distinction?

%% Q59
Q's percentages: I $=50\%$, II $=63/70=90\%$, III $=82\%$, IV $=60\%$. Distinction ($\ge75\%$) in \textbf{II and III}.

%% Q60 q
Approximate \% results of candidates P and Q will respectively, be:

%% Q60
Both candidates total $221/310$: P $=48+56+75+42=221$ and Q $=40+63+82+36=221$, so each scores about $71.3\%$. Among the options, the key takes \textbf{71\% and 70\%} as the approximate results.
